\section{Stienstra coefficients and Artin--Mazur height}
\label{sec:classical-height}
Let $p$ be any prime, including $2$, and let $F$ be a homogeneous
polynomial of degree $d$ in $d\ge3$ variables over $\mathbb F_p$,
defining a smooth projective hypersurface $X$.
Choose a homogeneous lift $\widetilde F\in\mathbb Z_p[x_1,\ldots,x_d]$
and put $z=x_1\cdots x_d$. Define
\begin{equation}
 a_m=a_m(\widetilde F)=[z^{m-1}]\widetilde F^{\,m-1}\quad(m\ge1).
 \label{eq:H1}
\end{equation}
Thus $a_1=1$ and $a_p\bmod p=h_p(F)$.

\subsection{The source theorem and its hypotheses}
Stienstra's Theorem~1, printed p.~907, with its base-change statement
on p.~908 \citep{Stienstra}, gives an integral one-dimensional formal
group law for the Artin--Mazur group of $X$ with logarithm
\begin{equation}
 \ell(T)=\sum_{m\ge1}\frac{a_m}{m}T^m.                 \label{eq:H2}
\end{equation}
His general theorem concerns a homogeneous regular sequence
$F_1,\ldots,F_r$ over a noetherian ring $K$, with flat projective
intersection, and degrees $d_i\ge \sum_i d_i-N\ge1$ in $\mathbb P^N$.
The logarithm statement additionally assumes that $K$ is flat over
$\mathbb Z$. In our specialization $r=1$, $N=d-1$, and the index set in
his theorem is the singleton $(1,\ldots,1)$, which gives exactly \eqref{eq:H1}.
A primitive hypersurface equation over $\mathbb Z_p$ gives a flat
hypersurface, so these hypotheses hold for any lift of the nonzero
equation $F$.

Artin--Mazur prove representability for projective complete intersections
of dimension at least two in the special case following
II, Corollary~(4.2), pp.~109--110; II, Corollary~(4.3), p.~110 identifies
the Dieudonn\'e module with Witt-vector cohomology
\citep{AM}. No crystalline comparison is needed for the
height-one or height-two computation below.

\subsection{An elementary logarithm lemma}
Vlasenko's general congruence theorem already determines height from
these valuations: in her notation $b_{p^j-1}=a_{p^j}$, finite height
$h$ gives $v_p(a_{p^j})\ge j-\lfloor j/h\rfloor$, with equality when
$h\mid j$, and infinite height gives $v_p(a_{p^j})\ge j$
\citep[Theorem~2(i), arXiv version~5, p.~3]{Vlasenko2019}.
The following proof also identifies the leading multiplication-by-$p$
coefficient in the chosen coordinate, which fixes the quotient used
in the counting interpretation.
\begin{lemma}
\label{lem:logarithm-height}
Let $\mathcal F$ be a one-dimensional commutative formal group law over
$\mathbb Z_p$ whose logarithm has the form
$\ell(T)=\sum_{m\ge1}a_mT^m/m$, with $a_m\in\mathbb Z_p$ and $a_1=1$.
Let $h$ be the height of its reduction, and write
$B(T)=[p]_{\mathcal F}(T)=\sum_{j\ge1}b_jT^j$.
For every $r\ge1$ for which $h\ge r$,
\begin{equation}
 a_{p^r}\in p^{r-1}\mathbb Z_p,\qquad
 b_{p^r}\equiv\frac{a_{p^r}}{p^{r-1}}\pmod p.         \label{eq:H3}
\end{equation}
Consequently, under that hypothesis,
\begin{equation}
 h=r\quad\Longleftrightarrow\quad
 \frac{a_{p^r}}{p^{r-1}}\not\equiv0\pmod p.           \label{eq:H4}
\end{equation}
\end{lemma}
\begin{proof}
We first recall the elementary power-of-$p$ property, including its
reason. If $g$ is a homomorphism of one-dimensional formal group laws
$\mathcal A,\mathcal B$ over a field of characteristic $p$, differentiation
of the homomorphism identity in its second variable at zero gives
\[
 g'(T)\,\partial_2\mathcal A(T,0)
 =\partial_2\mathcal B(g(T),0)\,g'(0).
\]
Both displayed partial derivatives are units. Thus $g'(0)=0$ implies
$g'(T)=0$ and $g(T)=g_1(T^p)$. The original identity and the injectivity
of substitution $(U,V)\mapsto(U^p,V^p)$ show that $g_1$ is a
homomorphism from the Frobenius twist of $\mathcal A$ to $\mathcal B$.
Iteration proves that the first nonzero term of a nonzero $p$-series
has degree $p^h$. It also proves that if its coefficient at $p^r$
vanishes after $r$ such factorizations, one can factor once more.
The zero $p$-series has height infinity.

Suppose $h\ge r$, and set $M=p^r$. Then $b_j\in p\mathbb Z_p$ for
$1\le j<M$. Compare the coefficient of $T^M$ in
$\ell(B(T))=p\ell(T)$:
\begin{equation}
 b_M+\sum_{j=2}^{M}\frac{a_j}{j}[T^M]B(T)^j
       =\frac{a_M}{p^{r-1}}.                        \label{eq:H5}
\end{equation}
Every term contributing to $[T^M]B^j$ uses $j$ coefficients $b_i$
with $i\le M-j+1<M$, and is therefore divisible by $p^j$.
Since $j-v_p(j)\ge1$ for $j\ge2$, each summand in \eqref{eq:H5} is an
integral multiple of $p$. This proves both the integrality assertion
and \eqref{eq:H3}. The power-of-$p$ property then gives \eqref{eq:H4}.
There is no exception at $p=2$.
\end{proof}

Applied to \eqref{eq:H2}, this proves the two tests
\begin{equation}
 h(X)=1\iff a_p\not\equiv0\pmod p,                  \label{eq:H6}
\end{equation}
and, on the locus $a_p\equiv0\pmod p$,
\begin{equation}
 h(X)=2\iff a_{p^2}/p\not\equiv0\pmod p.            \label{eq:H7}
\end{equation}
Here $a_{p^2}$ means the coefficient in $\widetilde F^{p^2-1}$,
not the square of $a_p$. Formula \eqref{eq:H7} requires that coefficient
modulo $p^2$.

\subsection{Lift independence and a universal second equation}
For a degree-$d$ polynomial $A$ in $d$ variables, write
$A_1(A)=[z^{p-1}]A^{p-1}$ and
$A_2(A)=[z^{p^2-1}]A^{p^2-1}$.
In characteristic $p$,
\begin{equation}
 A_2(A)=A_1(A)^{p+1}.                              \label{eq:H8}
\end{equation}
Indeed, factor $A^{p^2-1}=A^{p-1}(A^{p-1})^p$.
A contributing exponent in the first factor is congruent to $p-1$
in every coordinate. Each coordinate is at least $p-1$, while the
total degree is $d(p-1)$, forcing every coordinate to equal $p-1$.

Consequently, if the coefficients of $A$ are independent integer
variables, the polynomial
\begin{equation}
 H_2(A)=\frac{A_2(A)-A_1(A)^{p+1}}p                \label{eq:H9}
\end{equation}
has integer coefficients. On $H_1=A_1\bmod p=0$, its reduction agrees
with $a_{p^2}/p\bmod p$, since $a_p^{p+1}/p$ is divisible by $p$.
Thus height at least three is cut out, on the smooth parameter locus,
by $H_1=H_2=0$, and height exactly two by $H_1=0$, $H_2\ne0$.
This states equations for the loci; it makes no assertion about their
codimension, reducedness, smoothness, or nonemptiness.

There is also a direct check of the quotient's independence of the
chosen lift. Write $\widetilde F'=\widetilde F+pH$. Modulo $p^2$,
\[
 a_{p^2}(\widetilde F')-a_{p^2}(\widetilde F)
 \equiv p(p^2-1)[z^{p^2-1}]H\widetilde F^{p^2-2}.
\]
The same exponent argument, now applied to
$HF^{p-2}(F^{p-1})^p$, gives
\begin{equation}
 \frac{a_{p^2}(\widetilde F')-a_{p^2}(\widetilde F)}p
 \equiv
 -h_p(F)^p\,[z^{p-1}]HF^{p-2}\pmod p.             \label{eq:H10}
\end{equation}
It vanishes when $h_p(F)=0$. This proof includes $p=2$, where
$F^{p-2}=1$. Hence \eqref{eq:H7} is an intrinsic test of the special fibre,
not a choice of lift.

\subsection{A precise $p$-typical recurrence}
For clarity, one can obtain a $p$-typical logarithm without assuming
that the original logarithm \eqref{eq:H2} already is $p$-typical.
For a prime $q\ne p$ and an integral curve $\gamma(T)$ of $\mathcal F$,
apply the operator
\begin{equation}
 \gamma(T)\longmapsto
 \gamma(T)-_{\mathcal F}[1/q]_{\mathcal F}
       \Bigl(\mathop{\sum\nolimits^{\mathcal F}}_{\zeta^q=1}
                         \gamma(\zeta T)\Bigr).    \label{eq:H11}
\end{equation}
The sum can be formed over a finite unramified extension containing
the $q$th roots of unity; it is invariant and descends.
The inverse of $[q]_{\mathcal F}$ is integral because its tangent
coefficient $q$ is a unit. On logarithms, \eqref{eq:H11} removes precisely
the degrees divisible by $q$. Starting with $\gamma(T)=T$ and taking
the product over all primes $q\ne p$ is well defined degree by degree,
and gives an integral strict coordinate change $c(T)=T+O(T^2)$.
The logarithm in that coordinate is
\begin{equation}
 L(T)=\sum_{r\ge0}\lambda_rT^{p^r},\qquad
 \lambda_0=1,\quad \lambda_r=a_{p^r}/p^r.          \label{eq:H12}
\end{equation}

Let $\mathcal F_p$ be this formal group law. There are unique integral
parameters $w_r$, with $w_0=p$, such that
\begin{equation}
 [p]_{\mathcal F_p}(T)
 =\mathop{\sum\nolimits^{\mathcal F_p}}_{r\ge0}w_rT^{p^r}.
                                                               \label{eq:H13}
\end{equation}
To construct them, successively subtract the leading term of the
remaining curve using formal-group subtraction. Its logarithm has only
$p$-power degrees, so after the degree $p^r$ term is removed, the next
possible leading degree is $p^{r+1}$. The removed leading coefficient
is integral, which proves the assertion.

Taking logarithms of \eqref{eq:H13} gives the exact recurrence
\begin{equation}
 p\lambda_n=\sum_{i=0}^n\lambda_{n-i}w_i^{p^{n-i}},
 \qquad
 w_n=(p-p^{p^n})\lambda_n
        -\sum_{i=1}^{n-1}\lambda_{n-i}w_i^{p^{n-i}}.              \label{eq:H14}
\end{equation}
In particular,
\begin{equation}
 w_1=(1-p^{p-1})a_p,\qquad
 w_2=(1-p^{p^2-1})\frac{a_{p^2}}p
                   -\frac{a_p}{p}w_1^p.                       \label{eq:H15}
\end{equation}
The sum in \eqref{eq:H13} is formal-group addition, not ordinary addition.
Dropping the $p^{p^n}\lambda_n$ term while retaining the same
parameters would give an incorrect exact recurrence.
Modulo $p$, $w_1=H_1$ and $w_2=H_2$; the latter follows from
$(1-p^{p-1})^p\equiv1\pmod{p^2}$.
On $H_1=0$, \eqref{eq:H15} reduces to \eqref{eq:H7}.
The direct proof \eqref{eq:H5} establishes the height tests without requiring
this choice of coordinate.

\subsection{Finite tests and the crystalline hypothesis boundary}
For a prescribed cutoff $R$, formulas \eqref{eq:H3}--\eqref{eq:H4} give a finite procedure:
test $a_p\bmod p$, then, on its zero locus, test $a_{p^2}\bmod p^2$
after division by $p$, and continue through
$a_{p^R}/p^{R-1}\bmod p$.
The result is the exact height if it is at most $R$, or the statement
$h>R$. A finite cutoff alone does not identify infinite height.

For the smooth projective complete intersections here, one can choose
a smooth proper lift to $\mathbb Z_p$. Indeed, a homogeneous lift of
the nonzero equation gives a flat projective hypersurface. Its special
fibre is smooth, and properness makes the image of the closed
nonsmooth locus closed in $\operatorname{Spec}\mathbb Z_p$.
That image misses the closed point and must therefore be empty.
Artin--Mazur III,
Corollary~(1.9)(a), pp.~114--115, and III, Corollary~3.3 with
Remark~2, pp.~118--119, verify the hypotheses needed to identify the
Dieudonn\'e module, after inverting $p$, with the slope-$<1$ part of middle crystalline
cohomology \citep{AM}. In finite height its rank is the
height, so the middle crystalline rank supplies a finite upper bound.
Invoking that comparison for a general variety without the required
representability and lifting hypotheses would not be justified by
this citation.
